%======================================================================
%  Cambios al documento por las decisiones A1 y A2, y ajustes de
%  coherencia con el prototipo construido.
%
%  Cada bloque trae, en comentarios, el capítulo y la sección donde va,
%  el párrafo que reemplaza y el motivo. Se copia el texto que sigue al
%  comentario en el lugar indicado del proyecto LaTeX.
%======================================================================


%----------------------------------------------------------------------
%  CAMBIO 1 (decisión A1)
%  Capítulo 3, sección "Alcances Proyecto de BD".
%  Reemplaza el párrafo completo que empieza "El proyecto comprende el
%  análisis de toda la empresa...".
%  Motivo: el párrafo original dice que no se desarrollan las
%  aplicaciones de los usuarios. El prototipo incluye una aplicación web
%  mínima que demuestra los casos de uso del módulo; solo las
%  aplicaciones completas (por ejemplo, la móvil con funcionamiento sin
%  conexión) quedan fuera del alcance.
%----------------------------------------------------------------------

El proyecto comprende el análisis de toda la empresa, su diseño
conceptual global y, para el módulo priorizado, el diseño lógico, el
diseño físico y un prototipo funcional con tablas, restricciones,
consultas e informes, acompañado de una aplicación web mínima. Los demás
módulos conservan su análisis y su diseño conceptual
(anexos~\ref{anx:analisis} y~\ref{anx:global}) para implementarse
después. La aplicación web demuestra los casos de uso del módulo sobre la
base de datos: es una sola aplicación, con una sección para cada grupo de
usuarios (mesa de servicio, centro de distribución, mensajeros,
gerencia y clientes) y un portal público de rastreo por número de guía, y
se conecta a la base de datos mediante JDBC. Con ella el proyecto
integra la base de datos a una aplicación y diseña una aplicación
soportada en un motor relacional, dos propósitos del programa del curso.
Las aplicaciones completas, como la aplicación móvil del mensajero con
funcionamiento sin conexión, quedan fuera del alcance.


%----------------------------------------------------------------------
%  CAMBIO 2 (decisión A2)
%  Capítulo 5, sección "Conversión del Diseño Lógico al Diseño Físico",
%  subsección "Modelo Relacional al Modelo Computacional (DBMS)".
%  Reemplaza el párrafo completo "Usuarios, roles y permisos" (empieza
%  "\textbf{Usuarios, roles y permisos.} Cada grupo de usuarios del
%  módulo tiene un rol...").
%  Motivo: el párrafo original dice que los usuarios internos tienen una
%  cuenta propia en la base de datos. Ahora el servidor de aplicaciones
%  se conecta con una cuenta de servicio por rol y fija en cada conexión
%  el usuario real como identificador de cliente de la sesión de Oracle.
%  Se conserva lo demás: roles y permisos, RN16, esquema propietario y
%  tareas del administrador.
%----------------------------------------------------------------------

\textbf{Usuarios, roles y permisos.} Cada grupo de usuarios del módulo
tiene un rol con los permisos de la tabla~\ref{tab:permisos} traducidos a
privilegios: crear es insertar, actualizar es modificar y leer es
consultar. Ningún rol tiene permiso de borrar (RN16). Las tablas
pertenecen a un esquema propietario, BD\_ENLACE\_EXPRESS, con el que no
trabaja ningún usuario. El administrador de la base de datos crea los objetos,
carga los catálogos y, mientras se integran los módulos M2 a M4,
registra y actualiza los datos de referencia (clientes, sedes,
vehículos y empleados) a solicitud de las áreas dueñas
(tabla~\ref{tab:interfaces}). Los usuarios no tienen una cuenta propia
en la base de datos. El servidor de aplicaciones se conecta con una
cuenta de servicio por rol (cinco cuentas, como svc\_mesa y svc\_centro,
cada una con solo su rol) y elige la cuenta según el rol del usuario que
inició sesión en la aplicación; las consultas públicas del portal, que se
hacen antes de iniciar sesión, usan la cuenta de menor privilegio, la del
rol cliente. En cada conexión, el servidor fija además el identificador
de cliente de la sesión de Oracle con el usuario real, de modo que la
base de datos sabe quién hizo cada operación y puede auditarla. Así no
hay que crear ni administrar en la base de datos una cuenta por cada
empleado y cliente: el control por usuario lo hace la aplicación (por
ejemplo, el mensajero solo ve sus despachos y el cliente solo sus
envíos), y Oracle sigue limitando, mediante la cuenta de servicio, lo que
cada rol puede hacer.


%======================================================================
%  AJUSTES DE COHERENCIA CON EL PROTOTIPO CONSTRUIDO
%======================================================================

%----------------------------------------------------------------------
%  CAMBIO 3
%  Capítulo 6, introducción (antes de la sección "Algebra relacional de
%  consultas MySQL").
%  Reemplaza el párrafo que empieza "Este capítulo define cómo se usará
%  la base de datos del módulo...".
%  Motivo: dice que del prototipo solo se define la arquitectura y que su
%  construcción se desarrollará en la siguiente entrega. Ahora el
%  capítulo incluye la construcción del prototipo, la creación de la base
%  de datos y los resultados.
%----------------------------------------------------------------------

Este capítulo define cómo se usará la base de datos del módulo de
operación de envíos: expresa en álgebra relacional y en SQL las
consultas que necesitan sus usuarios, y especifica las tablas,
relaciones, transacciones e informes. Luego presenta el prototipo
funcional (sección~\ref{sec:prototipo}): su arquitectura, la creación de
la base de datos del módulo en el gestor, los datos de prueba, los
resultados de las consultas, las transacciones y los informes, y la
aplicación web que demuestra los casos de uso.


%----------------------------------------------------------------------
%  CAMBIO 4
%  Capítulo 6, sección "Prototipo funcional de la Base de Datos",
%  subsección "Arquitectura del prototipo" (último párrafo de la
%  subsección, después de la mención de ORDS).
%  Reemplaza el párrafo que empieza "El resto del prototipo se
%  desarrollará en la siguiente entrega...".
%  Motivo: la construcción ya no queda para otra entrega; el párrafo
%  pasa a ser el puente hacia las subsecciones siguientes del prototipo
%  (si la sección 6.3 nueva ya abre con su propia introducción, este
%  párrafo puede eliminarse).
%----------------------------------------------------------------------

Las subsecciones siguientes construyen el prototipo con esta
arquitectura: se crea la base de datos del módulo en el gestor, con los
objetos físicos del capítulo~5; se cargan datos de prueba; y se ejecutan
las consultas de la sección~\ref{sec:algebra}, las transacciones de la
tabla~\ref{tab:transacciones} y los informes de la
tabla~\ref{tab:informes}, y se muestran sus resultados.


%----------------------------------------------------------------------
%  CAMBIO 5 (opcional)
%  Capítulo 5, subsección "Recursos Hardware, Software, Aplicaciones y
%  BD, Comunicaciones, Ingeniería Web", párrafo "Aplicaciones y base de
%  datos".
%  Reemplaza la última oración del párrafo: "Las aplicaciones completas
%  no forman parte del alcance del proyecto (sección~\ref{sec:modulo});
%  el prototipo funcional (sección~\ref{sec:prototipo}) implementa sus
%  funciones principales sobre la base de datos del módulo."
%  Motivo: la tabla lista cuatro aplicaciones, y el prototipo las
%  demuestra en una sola aplicación web con secciones por grupo de
%  usuarios (decisión A1). La oración original no se contradice, pero
%  con este ajuste queda explícito. El resto del párrafo ("se conecta con
%  el rol de ese grupo") sigue siendo cierto con las cuentas de servicio
%  de A2.
%----------------------------------------------------------------------

Las aplicaciones completas no forman parte del alcance del proyecto
(sección~\ref{sec:modulo}); el prototipo funcional
(sección~\ref{sec:prototipo}) demuestra sus funciones principales sobre
la base de datos del módulo en una sola aplicación web, con una sección
para cada grupo de usuarios y el portal público de rastreo.
